% outline.md, lines 9--35: one frame with five successive overlays.
\begin{frame}{Classical Verification of Quantum Computation}
  \centering
  \begin{tikzpicture}[
      >={Latex[length=3mm]},
      party/.style={circle,draw=CVQCInk,fill=CVQCFog,very thick,
        minimum size=2.5cm,text width=2.1cm,inner sep=0pt,
        align=center,font=\large\bfseries,text=CVQCInk},
      callout/.style={fill=CVQCPale,rounded corners=4mm,
        minimum width=4.7cm,minimum height=0.95cm,
        align=center,text=CVQCInk,inner sep=3mm},
      communication/.style={->,line width=1.5pt,draw=CVQCTeal}
    ]
    % Reserve the same space on every overlay so the parties never move.
    \path[use as bounding box] (-6.5,-2.85) rectangle (6.5,3.1);

    \node[party] (verifier) at (-4.15,0) {Verifier};
    \node[party,fill=CVQCPale] (prover) at (4.15,0) {Prover};
    \quantumpartymark{prover}

    % Pale bubbles and angled tails follow Mending-Slides/diagrams/results.tex.
    % Each annotation belongs only to its own stage of the interaction.
    % Input and output share a raised bubble, giving its tail a gentler angle.
    \only<2,4-5>{
      \node[callout] (annotation) at (-3.15,2.45) {};
      \path[fill=CVQCPale]
        ($(annotation.south)+(-0.65,0.02)$) --
        ($(verifier.north)+(-0.15,0.06)$) --
        ($(annotation.south)+(-0.05,0.02)$) -- cycle;
      \node[font=\large,text=CVQCInk] at (annotation.center) {%
        \only<2>{$C(x)$: Quantum circuit}%
        \only<4>{$y$: circuit output}%
        \only<5>{$y \gets C(x)\;\text{or}\;\textcolor{CVQCCoral}{\mathsf{reject}}$}%
      };
    }

    \only<3>{
      % Leave a gap to both circles, even at their widest point.
      \draw[communication] (-2.35, 1.10) -- ( 2.35, 1.10);
      \draw[communication] ( 2.35, 0.50) -- (-2.35, 0.50);
      \draw[communication] (-2.35,-0.10) -- ( 2.35,-0.10);
      \node[font=\large,text=CVQCTeal] at (0,-0.60) {$\vdots$};
      \draw[communication] ( 2.35,-1.10) -- (-2.35,-1.10);
    }

  \end{tikzpicture}

  \note<1>{
    \footnotesize
    Think of a classical client using a remote quantum server. The server can
    perform computations beyond the client's reach, but how can the client
    check that the computation was performed correctly? This is also a natural
    question when demonstrating quantum advantage: the very separation that
    makes the computation interesting makes its output difficult to verify.
  }
  \note<2>{
    \footnotesize
    The client has a description of a quantum circuit $C$ and a classical input
    $x$. Here $C(x)$ denotes running that circuit and measuring its output to
    obtain a classical sample. Knowing the circuit does not mean the client can
    efficiently execute it.
  }
  \note<3>{
    \footnotesize
    The verifier and prover exchange classical messages. The alternating
    arrows and ellipsis stand for an interactive protocol, potentially with
    many rounds; the arrows are schematic, not an exact message count.
    The verifier remains entirely classical; the prover does the quantum work.
  }
  \note<4>{
    \footnotesize
    At the end, the verifier produces a classical output $y$. The goal is to
    verify this output through the interaction, without having to reproduce
    the quantum computation locally.
  }
  \note<5>{
    \footnotesize
    This is shorthand for the desired sample-or-reject behavior, up to the
    protocol's error. Completeness: an honest prover makes the verifier accept
    and return a sample from $C(x)$ with high probability. Soundness: even an
    efficient cheating prover cannot make the verifier's output distinguishable
    from an ideal experiment that independently samples from $C(x)$ or rejects.
    The ideal experiment has the same acceptance probability as the real one.
    The guarantee concerns the entire output distribution, including rejection;
    a single sample is not individually classified as correct or incorrect.
  }
\end{frame}
